\documentclass[12pt,a4paper,english
% ,twoside,openright
]{tutthesis}
%\documentclass[12pt,a4paper,finnish]{tutthesis}

% Note that you must choose either Finnish or English here and there in this
% file.
% Other options for document class
  % ,twoside,openright   % If printing on both sides (>80 pages)
  % ,twocolumn           % Can be used in lab reports, not in theses

% Ensure the correct Pdf size (not needed in all environments)
\special{papersize=210mm,297mm}


% LaTeX file for BSC/MSc theses and lab reports.
% Requires the class file (=template) tutthesis.cls and figure files,
% either tut-logo, exampleFig (as pdf or eps) and example_code.c
% Author: Sami Paavilainen (2006)
% Modified: Heikki Huttunen (heikki.huttunen@tut.fi) 31.7.2012.
%           Erno Salminen, @tut.fi, 2014-08-15
%             - added text snippets from the writing guide
%             - added lots of comments: both tips and alternative styles
%             - added an example table
%             - and so on...

%
% Define your basic information
%
\author{Jaakko Pasanen}
\title{NLP for Customer Support Agent} % primary title (for front page)
\titleB{Luonnollisen kielen ymmärrys asiakapalveluagentilla}     % translated title for abstract
\thesistype{Master of Science thesis} % or Bachelor of Science, Laboratory Report... 
\examiner{Ari Visa} % without title Prof., Dr., MSc or such

% Put your thesis' main language last
% http://mirrors.ctan.org/macros/latex/required/babel/base/babel.pdf
\usepackage[finnish, main=english]{babel}

% http://www.ctan.org/pkg/biblatex
\usepackage[
  style=authoryear,
  maxcitenames=2,
  backend=biber,
  firstinits=true
]{biblatex}
\bibliography{thesis_refs.bib}
%% Note that option style=numeric works as well

\usepackage{fontspec}


% You can also add your own commands
\newcommand\todo[1]{{\color{red}!!!TODO: #1}} % Remark text in braces appears in red
\newcommand{\angs}{\textsl{\AA}}              % , e.g. slanted symbol for Ångstöm

% Preparatory content ends here


\pagenumbering{roman} % was: {Roman}
\pagestyle{headings}
\begin{document}

% Special trick so that internal macros (denoted with @ in their name)
% can be used outside the cls file (e.g. \@author)
\makeatletter

%
% Create the title page.
% First the logo. Check its language.
\thispagestyle{empty}
\vspace*{-.5cm}\noindent
\includegraphics[width=8cm]{tty_tut_logo}   % Bilingual logo

% Then lay out the author, title and type to the center of page.
\vspace{6.8cm}
\maketitle
\vspace{7.7cm} % -> 6.7cm if thesis title needs two lines

% Last some additional info to the bottom-right corner
\begin{flushright}  
  \begin{minipage}[c]{6.8cm}
    \begin{spacing}{1.0}
      %\textsf{Tarkastaja: Prof. \@examiner}\\
      %\textsf{Tarkastaja ja aihe hyväksytty}\\ 
      %\textsf{xxxxxxx tiedekuntaneuvoston}\\
      %\textsf{kokouksessa dd.mm.yyyy}\\
      \textsf{Examiner: Prof. \@examiner}\\
      \textsf{Examiner and topic approved by the}\\ 
      \textsf{Faculty Council of the Faculty of}\\
      \textsf{xxxx}\\
      \textsf{on 30th July 2014}\\
    \end{spacing}
  \end{minipage}
\end{flushright}

% Leave the backside of title page empty in twoside mode
\if@twoside
\clearpage
\fi

%
% Use Roman numbering I,II,III... for the first pages (abstract, TOC,
% termlist etc)
\pagenumbering{Roman} 
\setcounter{page}{0} % Start numbering from zero because command 'chapter*' does page break

% Some fields in abstract are automated, namely those with \@ (author,
% title in the main language, thesis type, examiner).
\chapter*{Abstract}

\begin{spacing}{1.0}
         {\bf \textsf{\MakeUppercase{\@author}}}: \@title\\   % use \@titleB when thesis is in Finnish
         \textsf{Tampere University of Technology}\\
         \textsf{\@thesistype, xx pages, x Appendix pages} \\
         \textsf{xxxxxx 201x}\\
         \textsf{Master's Degree Programme in xxx Technology}\\
         \textsf{Major: }\\
         \textsf{Examiner: Prof. \@examiner}\\ % 
         \textsf{Keywords: }\\
\end{spacing}


The abstract is a concise 1-page description of the work: what was the
problem, what was done, and what are the results. Do not include
charts or tables in the abstract.

Put the abstract in the primary language of your thesis first and then
the translation (when that is needed).


% Foreign students do not need Fininsh abstract (tiivistelmä). Move
% this before English abstract if thesis is in Finnish. Move also the
% otherlanguage command to the English abstract (if needed).

\begin{otherlanguage}{finnish} %  Following text in in 2nd language
\chapter*{Tiivistelmä} % Asterisk * turns numbering off

\begin{spacing}{1.0}
         {\bf \textsf{\MakeUppercase{\@author}}}: \@titleB\\  % or use \@title when thesis is in Finnish
         \textsf{Tampereen teknillinen yliopisto}\\
         \textsf{Diplomityö, xx sivua, x liitesivua}\\ %
         \textsf{xxxkuu 201x}\\
         \textsf{xxx koulutusohjelma}\\
         \textsf{Pääaine: }\\
         \textsf{Tarkastajat:  Prof. \@examiner}\\ % automated, if just 1 examiner
         \textsf{Avainsanat: }\\
\end{spacing}

The abstract in Finnish. Foreign students do not need this page.

Suomenkieliseen diplomityöhön kirjoitetaan tiivistelmä sekä suomeksi
että englanniksi.

Kandidaatintyön tiivistelmä kirjoitetaan ainoastaan kerran, samalla
kielellä kuin työ. Kuitenkin myös suomenkielisillä kandidaatintöillä
pitää olla englanninkielinen otsikko arkistointia varten.

\end{otherlanguage} % End on 2nd language part


\chapter*{Preface}

This document template conforms to Guide to Writing a Thesis at
Tampere University of Technology (2014) and is based on the previous
template. The main purpose is to show how the theses are formatted
using LaTeX (or \LaTeX ~ to be extra fancy) .

The thesis text is written into file \texttt{d\_tyo.tex}, whereas
\texttt{tutthesis.cls} contains the formatting instructions. Both
files include lots of comments (start with \%) that should help in
using LaTeX. TUT specific formatting is done by additional settings on
top of the original \texttt{report.cls} class file. This example needs
few additional files: TUT logo, example figure, example code, as well
as example bibliography and its formatting (\texttt{.bst}) An example
makefile is provided for those preferring command line. You are
encouraged to comment your work and to keep the length of lines
moderate, e.g. <80 characters. In Emacs, you can use \texttt{Alt-Q} to
break long lines in a paragraph and \texttt{Tab} to indent commands
(e.g. inside figure and table environments). Moreover, tex files are
well suited for versioning systems, such as Subversion or Git.  
% \url{http://www.ctan.org/tex-archive/info/lshort/english/lshort.pdf}

Acknowledgements to those who contributed to the thesis are generally
presented in the preface. It is not appropriate to criticize anyone in
the preface, even though the preface will not affect your grade. The
preface must fit on one page. Add the date, after which you have not
made any revisions to the text, at the end of the preface.

~ 
% Tilde ~ makes an non-breakable spce in LaTeX. Here it is used to get
% two consecutive paragraph breaks

Tampere, 11.8.2014

~

On behalf of the working group, Erno Salminen


% Add the table of contents, optioanlly also the lists of figures,
% tables and codes.

%\renewcommand\contentsname{Sisällys} % Set Finnish name, remove this if using English
\setcounter{tocdepth}{3}              % How many header level are included
\tableofcontents                      % Create TOC


%
% Term and symbol exaplanations use a special list type
%
\chapter*{List of abbreviations and symbols}
\markboth{}{}                                % no headers
%\chapter*{Lyhenteet ja merkinnät}

% You don't have to align these with whitespaces, but it makes the
% .tex file more readable
\begin{termlist}
\item[NER] Named entity recognition
\item[NLP] Natural Language Processing
\item[POS] Part-of-speech also called lexical category
\item[TUT] Tampere University of Technology
\end{termlist}


\chapter*{Terms and Notes}
\begin{termlist}
\item[Bag-of-words]
\item[Constituent] In syntactic analysis, a constituent is a word or a group of words that function(s) as a single unit within a hierarchical structure. Many constituents are phrases. \textit{Yesterday I saw \textbf{an orange bird with a white neck}}
\item[Lemmatisation] Process of finding the base form of a word, e.g. flew -> fly
\item[Parsing] Within computational linguistics the term is used to refer to the formal analysis by a computer of a sentence or other string of words into its constituents, resulting in a parse tree showing their syntactic relation to each other, which may also contain semantic and other information.
\item[POS-tagging] Process of marking up a word to particular part-of-speech (nouns, verbs, etc...) based on both its definition and its context.
\item[Skip-gram]
\item[Tree bank] Parsed text corpus that annotates syntactic or semantic sentence structure. Contains trees for sentences where phrases in a sentence are structured in a tree of syntactic or semantic relations. Very useful for training POS-taggers etc...
\item[Word vector] vector('Paris') - vector('France') + Vector('Italy') -> vector('Rome')
\end{termlist}


% The actual text begins here and page numbering changes to 1,2...
% Leave the backside of title empty in twoside mode
\if@twoside
%\newpage
\cleardoublepage
\fi

\pagenumbering{arabic}
\setcounter{page}{1} % Start numbering from zero because command
                     % 'chapter*' does page break
\renewcommand{\chaptername}{} % This disables the prefix 'Chapter' or
                              % 'Luku' in page headers (in 'twoside'
                              % mode)


\chapter{Introduction}
\label{ch:intro} 
Testing citation \cite{Andor2016}

\chapter{Natural Language Processing}
\label{ch:natural_language_processing}

\section{POS Tagging}
\label{se:pos_tagging}


\section{Dependecy Parsing}
\label{se:dependecy_parsing}


\section{Co-Reference Parsing}
\label{se:co-reference_parsing}


\section{Sentence Segmentation}
\label{se:sentence_segmentation}


\section{Lemmatisation}
\label{se:lemmatisation}


%
% The bibliography, i.e the list of references (3 options available)
%
\newpage


% Extra for Finnish theses

\renewcommand{\bibname}{Bibliography}     % Bilingual babel puts Finnish ``Kirjallisuttaa'' otherwise. Strange...
%\renewcommand{\bibname}{Lähteet}         % Set Finnish header, remove this if using English
%\addcontentsline{toc}{chapter}{Lähteet}  % Include this in TOC
\addcontentsline{toc}{chapter}{\bibname}  % Include this in TOC


\printbibliography                  % a) heading in English
%\printbibliography[title=Lähteet]   % b) heading in Finnish
%\addtocontents{toc}{%               % b) add Finnish heading to table of contents
%\protect\noindent Lähteet\protect\par
%} 


%
% Appendices are optional. 
% This part is semi-ugly at the moment. Please give feedback if can
% improve it.
\appendix
\pagestyle{headings}
% \renewcommand{\appendixname}{Liite} % Extra. Set Finnish prefix for page header

%
% a) Not-so-handy way, but at least it works
% 
\def\appA{APPENDIX A. Something extra} % Define the name and numbering manually
\chapter*{\appA}                       % Create chapter heading
\markboth{\appA}{\appA}                % Set page header
\addcontentsline{toc}{chapter}{\appA}  % Include this in TOC
% Note that \label does not work with unnumbered chapter

Appendices are purely optional.  All appendices must be referred to in
the body text

\def\appB{APPENDIX B. Something completely different} % Define another new command
\chapter*{\appB}                       % As above, but use \appB instead of \appA
\label{app:B}
\markboth{\appB}{\appB}                     
\addcontentsline{toc}{chapter}{\appB}  


You can append to your thesis, for example, lengthy mathematical
derivations, an important algorithm in a programming language, input
and output listings, an extract of a standard relating to your thesis,
a user manual, empirical knowledge produced while preparing the
thesis, the results of a survey, lists, pictures, drawings, maps,
complex charts (conceptual schema, circuit diagrams, structure charts)
and so on.


%
% b) The other option is to use numbered chapter and our baseline
% template report.cls numbers them as A, B... The heading and TOC do
% not include prefix 'Appendix' although the page header does.
%\chapter{name of the appendix}
%\label{app:A}                          % For cross-references



\end{document}

